\section{Корпус и его характеристики}

Тематика корпуса --- \textbf{кино и сериалы}: фильмы, сериалы, мультфильмы и
кинематографические персоналии. Источника два:

\begin{itemize}
    \item \href{https://ru.wikipedia.org/}{\textbf{ru.wikipedia.org}} --- энциклопедические
          статьи о фильмах, сериалах и кинематографистах;
    \item \href{https://www.film.ru/}{\textbf{www.film.ru}} --- рецензии на фильмы и
          сериалы, тематические подборки и интервью.
\end{itemize}

Источники дополняют друг друга по жанру текста: энциклопедическая статья сообщает факты
о фильме, рецензия --- мнение о нём. Для поиска это разные запросы: \enquote{кто снял} и
\enquote{стоит ли смотреть}.

\subsection*{Из чего состоит текст}

\textbf{Википедия}, кодировка UTF-8. Статьи собраны по устойчивому шаблону, разделы почти
всегда одни и те же: \emph{Сюжет}, \emph{В ролях}, \emph{Съёмочная группа},
\emph{Создание}, \emph{Критика}, \emph{Награды}, \emph{Примечания}, \emph{Литература},
\emph{Ссылки}. Три последних раздела --- списки источников и внешних адресов, к
содержанию документа они отношения не имеют.

\textbf{film.ru}, кодировка UTF-8. Статья --- сплошной авторский текст из нескольких
абзацев без подзаголовков, с иллюстрациями между ними; в конце стоит подпись автора и
кнопки \enquote{Поделиться}. Подборки устроены как нумерованный список фильмов с
абзацем текста о каждом, интервью --- как чередование вопросов и ответов. Кроме статей
на сайте есть карточки фильмов, но текста в них --- один-два абзаца аннотации.

\subsection*{Разметка}

Обе площадки отдают обычный HTML, текст виден прямо в исходном коде страницы. Основной
текст лежит в одном контейнере: у Википедии это
\texttt{<div class="mw-content-ltr mw-parser-output">}, у film.ru ---
\texttt{<div class="wrapper\_articles\_text">} у статей и
\texttt{<div class="wrapper\_movies\_text">} у карточек фильмов. Всё, что вне
контейнера, --- меню, боковые панели, комментарии и подвал.

Внутри текста у Википедии встречаются карточка фильма
(\texttt{<table class="infobox">}) со структурированными полями, оглавление, сноски
\texttt{<sup class="reference">} и заголовки разделов \texttt{<h2>}. У film.ru внутри
контейнера только абзацы \texttt{<p>}, выделения \texttt{<strong>} и картинки.

\subsection*{Дополнительная мета-информация}

У Википедии: канонический адрес \texttt{<link rel="canonical">} и блок JS-конфигурации с
внутренним номером статьи \texttt{wgArticleId} и полным списком её категорий
\texttt{wgCategories} --- фактически готовая рубрикация. Даты документа в HTML нет.

У film.ru мета-информации больше. Кроме набора Open Graph (\texttt{og:title},
\texttt{og:description}, \texttt{og:url}, \texttt{og:image}) в странице есть разметка
schema.org в формате JSON-LD: у рецензии --- тип \texttt{Review} с автором, точной датой
публикации \texttt{datePublished}, оценкой фильма по десятибалльной шкале и описанием
самого фильма (\texttt{Movie}) с режиссёром и актёрами.

\subsection*{Почему не Кино-Театр.Ру}

Первым кандидатом на второй источник был \texttt{www.kino-teatr.ru}: карточки фильмов,
биографии актёров и авторские статьи. Пять его страниц сохранены в
\texttt{samples/kinoteatr}, и по ним видны две его особенности. Он отдаёт
страницы в \textbf{windows-1251} --- при чтении в UTF-8 весь текст превращается в мусор. И
содержательный текст там есть только у статей: биографии и карточки фильмов --- это
несколько предложений, от 8 до 46 слов.

От источника пришлось отказаться при обкачке во второй лабораторной работе: сайт закрыт
защитой ddos-guard, которая после примерно пятисот запросов отдаёт роботу капчу вместо
страниц и не снимает блокировку часами. На ручном скачивании пяти страниц это не
проявилось, поэтому второй источник стоит проверять серией запросов, а не одним.

\pagebreak
